\chapter{与AI对话的艺术：Prompting基础}

本章介绍结构化提示的基本方法，帮助读者写出高质量的初始 Prompt，并通过迭代持续改进 AI 输出。

\section{为什么“随便问”得不到好答案}

同一个问题使用不同的提问方式，可能得到截然不同的结果。模糊的提问会迫使 AI 猜测任务背景、目标和约束，而这些猜测很可能与真实需求不一致。

\section{Prompt的解剖学：一个好提示由什么组成}

\begin{enumerate}
  \item \textbf{角色（Role）}：告诉 AI 它应以什么身份完成任务。
  \item \textbf{上下文（Context）}：说明数据来源、字段含义和业务场景。
  \item \textbf{任务（Task）}：明确 AI 需要完成什么，任务应具体、可执行且有边界。
  \item \textbf{约束（Constraint）}：说明不能做什么，例如不要删除原始列、不要使用未来信息。
  \item \textbf{输出格式（Format）}：指定 Python 代码、Markdown 表格或逐步解释等输出形式。
\end{enumerate}

\section{提示工程的核心技巧}

\subsection{零样本与少样本提示}

根据任务复杂程度判断是否需要提供示例。当 AI 已经能够理解意图时，可以使用零样本提示；当任务规则特殊或输出格式严格时，应提供少量示例。

\subsection{思维链提示}

对于复杂任务，可以要求 AI 先列出分析思路，再生成代码或结论，以提升过程的可检查性。

\subsection{分步拆解}

将“帮我建立一个预测模型”拆成描述数据结构、提出特征工程建议、选择模型和验证结果等多个步骤。

\subsection{迭代修正}

基于 AI 的第一次输出编写修正 Prompt，指出具体问题、补充缺失约束，而不是每次重新从零开始。

\section{常见提示陷阱与反面案例集}

常见问题包括任务目标模糊、缺少必要上下文、未声明约束、一次要求完成过多任务，以及无法验证输出是否正确。

\section{实战练习}

选择一个真实任务，分别使用结构化提示、分步拆解和迭代修正完成该任务，并比较不同提示方式的结果。
